\section{Sémantika datalogu}
Zoberme si veľmi jednoduchú teóriu: 
\begin{verbatim}
    p = {(1)}.    
\end{verbatim}

Pre túto teóriu existuje niekoľko modelov (dokonca ich je nekonečne veľa). Najjednoduchší z nich je
$\inter \models_\true \verb|p|(1)$ a pre všetky ostatné formuly je interpretácia nepravdivá. Ďalšími
modelmi môžu byť napríklad:
\begin{itemize}
\item $\inter' \models_\true \verb|p|(i), \forall i \in \mathbb{N}$ a pre všetky ostatné formuly je
interpretácia nepravdivá,
\item $\inter'' \models_\true \verb|p|(i), \forall i \in \mathbb{R}$ a pre všetky ostatné formuly je
interpretácia nepravdivá,
\item $\inter''' \models_\true \verb|p|(1), \inter''' \models_\true \verb|p|(\text{Pes}), \inter'''
\models_\true p(47), \inter''' \models_\true \verb|p|(\text{Štvorec})$ a pre všetky ostatné formuly
je interpretácia nepravdivá.
\end{itemize}

Ako sémantiku našich datalogových programov budeme uvažovať \textit{well-founded model}.
Well-founded model je minimálny stabilný trojhodnotový model.

\subsubsection{}
\textit{Trojhodnotový} model znamená, že interpretácia priraďuje atomickým formulám hodnoty
\textit{true} $\true$, \textit{false} $\false$ a \textit{unknown} $\unknown$.

\subsubsection{}
\textit{Stabilný} model znamená, že zostane rovnaký po tzv. \textit{Gelfond-Lifschitz
Transformation}. Aplikácia GLT funguje nasledovne: zoberieme aktuálnu interpretáciu a pomocou nej z
programu odstránime tie pravidlá, ktorých negované podmienky sú podľa danej interpretácie pravdivé.
Z ostatných pravidiel potom odstránime negácie. Tým dostaneme program bez negácie, pre ktorý vieme
zostrojiť minimálny model. Ak sa tento minimálny model zhoduje s interpretáciou, z ktorej sme
vychádzali, hovoríme, že ide o stabilný model.

Zoberme si program
\begin{verbatim}
    a.
    b :- \+ c.
    c :- \+ b.
    d :- a, \+ c.
\end{verbatim}

a kandidátnu interpretáciu $\inter$, pre ktorú platí
$$\inter \models_\true \verb|a|, \inter \models_\true \verb|b|, \inter \models_\false \verb|c|,
\inter \models_\true \verb|d|\text{.}$$

Pri aplikácii GLT najprv odstránime tie pravidlá, ktorých negované podmienky sú podľa $\inter$
pravdivé. Pravidlo \verb|c :- \+ b.| sa preto odstráni, keďže \verb|b| je v interpretácii $\inter$
pravdivé. V ostatných pravidlách potom odstránime negácie, a dostaneme program
\begin{verbatim}
    a.
    b.
    d :- a.
\end{verbatim}

Minimálny model takto upraveného programu obsahuje práve atómy \verb|a|, \verb|b| a \verb|d|. Keďže
sa zhoduje s interpretáciou, z ktorej sme vychádzali, daná interpretácia je stabilný model pôvodného
programu.

\subsubsection{}
\textit{Minimálny} model znamená, že spomedzi všetkých modelov vyberáme ten, ktorý obsahuje najmenej
pravdivých atómov. Intuitívne teda nechceme odvodiť viac faktov, než je nutné. Ak existuje viac
modelov, uprednostníme ten najmenší.

Príklad minimálneho modelu pre teóriu
\begin{verbatim}
    p = {(1)}.
\end{verbatim}
je modelom napríklad interpretácia, v ktorej platí iba $p(1)$. Modelom je aj interpretácia, v ktorej
platí $p(1)$, $p(2)$ a $p(3)$. Minimálny model je však ten prvý, pretože obsahuje najmenej
pravdivých atómov.

\subsubsection{}
\textit{Well-founded} model môžeme chápať ako niečo, čo je spoločné pre všetky stabilné modely.
Atómy, ktoré sú pravdivé vo všetkých stabilných modeloch, označíme ako $\true$, atómy, ktoré sú
nepravdivé vo všetkých stabilných modeloch, označíme ako $\false$, a atómy, pri ktorých sa stabilné
modely rozchádzajú, ponecháme ako $\unknown$.

V nasledujúcom príklade ukážeme, ako sa dá nájsť well-founded model teórie. Nejde o formálny dôkaz.
Pre viacej o well-fouded model odporúčam prečítať \cite{WellFounded} a \cite{SemanticsQueries}. 

Budeme postupovať iteratívne. Začneme interpretáciou, v ktorej sú všetky
atómy nepravdivé, a budeme opakovane aplikovať GLT. Po každom kroku vezmeme minimálny model
získaného programu a ten použijeme ako novú interpretáciu. Keď sa interpretácia prestane meniť,
dostaneme výsledný model.

Príklad nájdenia well-founded modelu:
\begin{verbatim}
    a.
    b :- a.
    p :- \+ q.
    q :- \+ p.
    u :- v.
\end{verbatim}

Budeme uvažovať atómy \verb|a|, \verb|b|, \verb|p|, \verb|q|, \verb|u| a \verb|v|. Na začiatku
zoberieme interpretáciu $\inter_0$, v ktorej sú všetky tieto atómy nepravdivé.

Pri aplikácii GLT na $\inter_0$ sa neodstráni žiadne pravidlo, pretože v $\inter_0$ nie je pravdivý
ani \verb|p|, ani \verb|q|. Po odstránení negácií dostaneme program
\begin{verbatim}
    a.
    b :- a.
    p.
    q.
    u :- v.
\end{verbatim}
Jeho minimálny model je interpretácia $\inter_1$, v ktorej platí
$$\inter_1 \models_\true \verb|a|, \inter_1 \models_\true \verb|b|, \inter_1 \models_\true \verb|p|,
\inter_1 \models_\true \verb|q|, \inter_1 \models_\false \verb|u|, \inter_1 \models_\false \verb|v|\text{.}$$

Teraz aplikujeme GLT na interpretáciu $\inter_1$. Keďže v nej sú \verb|p| aj \verb|q| pravdivé,
pravidlá \verb|p :- \+ q.| aj \verb|q :- \+ p.| sa odstránia. Po odstránení negácií dostaneme program
\begin{verbatim}
    a.
    b :- a.
    u :- v.
\end{verbatim}
Jeho minimálny model je interpretácia $\inter_2$, v ktorej platí
$$\inter_2 \models_\true \verb|a|, \inter_2 \models_\true \verb|b|, \inter_2 \models_\false \verb|p|,
\inter_2 \models_\false \verb|q|, \inter_2 \models_\false \verb|u|, \inter_2 \models_\false \verb|v|\text{.}$$

Ak teraz ešte raz aplikujeme GLT na interpretáciu $\inter_2$, obe pravidlá s negáciou opäť zostanú,
a preto znovu dostaneme program
\begin{verbatim}
    a.
    b :- a.
    p.
    q.
    u :- v.
\end{verbatim}
Jeho minimálny model je opäť $\inter_1$. Dostali sme teda cyklus
$$\inter_0 \rightarrow \inter_1 \rightarrow \inter_2 \rightarrow \inter_3 \equiv \inter_1 \rightarrow \inter_2 \rightarrow \cdots$$

Vidíme, že atómy \verb|a| a \verb|b| sú pravdivé v každom kroku od $\inter_1$ ďalej, a preto im vo
well-founded modeli priradíme hodnotu $\true$. Atómy \verb|u| a \verb|v| nie sú pravdivé v žiadnom
kroku, a preto im priradíme hodnotu $\false$. Naopak, pri atómoch \verb|p| a \verb|q| sa interpretácie
striedajú: raz sú pravdivé a raz nepravdivé. Tým pádom ich nevieme jednoznačne určiť, a vo
well-founded modeli dostanú hodnotu $\unknown$.

Výsledný well-founded model tejto teórie je teda daný hodnotami
$$\verb|a| = \true, \qquad \verb|b| = \true, \qquad \verb|u| = \false, \qquad \verb|v| = \false, \qquad \verb|p| = \unknown, \qquad \verb|q| = \unknown\text{.}$$

Prehľadne to môžeme zapísať aj do tabuľky:

\begin{center}
\begin{tabular}{c|c|c|c|c|c}
Atóm & $\inter_0$ & $\inter_1$ & $\inter_2$ & $\inter_3$ & $\inter$ \\
\hline
\verb|a| & $\false$ & $\true$ & $\true$ & $\true$ & $\true$ \\
\verb|b| & $\false$ & $\true$ & $\true$ & $\true$ & $\true$ \\
\verb|u| & $\false$ & $\false$ & $\false$ & $\false$ & $\false$ \\
\verb|v| & $\false$ & $\false$ & $\false$ & $\false$ & $\false$ \\
\verb|p| & $\false$ & $\true$ & $\false$ & $\true$ & $\unknown$ \\
\verb|q| & $\false$ & $\true$ & $\false$ & $\true$ & $\unknown$ \\
\end{tabular}
\end{center}

